% =====================================================================
\chapter{Literature Review}
\label{ch:lit}
% =====================================================================

\section{Introduction}

Research on detecting machine-generated text began in earnest with the
release of GPT-2 in 2019 and has grown rapidly since the public release of
ChatGPT. Recent surveys~\cite{tang2024science,wu2025survey} organise the field
by the kind of evidence a detector uses and by the access it requires to the
generating model. For this project the literature and the market are reviewed
in seven areas:
\begin{enumerate}
  \item statistical and heuristic (zero-shot and feature-based) methods;
  \item supervised transformer-based classifiers;
  \item commercial detection products;
  \item watermarking and generation-time signals;
  \item benchmark datasets and the evidence they provide about generalisation;
  \item studies of robustness, reliability and fairness;
  \item efficient inference through cascades and early exit, which is the
        design principle that \veriscript{} borrows from outside the
        detection literature.
\end{enumerate}
Fig.~\ref{fig:timeline} places the main works on a timeline. Each area is
reviewed below, followed by a synthesis that identifies the research gap
addressed by this project.

\begin{figure}[H]
\centering
\begin{tikzpicture}[x=2.45cm,y=1cm]
  \draw[thick,->] (0,0) -- (5.6,0);
  \foreach \x/\yr in {0.3/2019,1.3/2020,2.3/2023,3.6/2024,4.9/2025}{
     \draw[thick] (\x,0.12) -- (\x,-0.12);
     \node[below,font=\small\bfseries] at (\x,-0.15) {\yr};}
  \node[blk,font=\scriptsize,text width=2.0cm,above] at (0.3,0.45)
     {GLTR~\cite{gehrmann2019gltr}\\GPT-2 detector~\cite{solaiman2019release}\\Grover~\cite{zellers2019grover}};
  \node[blk,font=\scriptsize,text width=1.9cm,below] at (1.3,-0.75)
     {Human vs.\ machine study~\cite{ippolito2020automatic}};
  \node[blk,font=\scriptsize,text width=2.2cm,above] at (2.3,0.45)
     {DetectGPT~\cite{mitchell2023detectgpt}\\Watermark~\cite{kirchenbauer2023watermark}\\HC3~\cite{guo2023hc3}\\DIPPER~\cite{krishna2023paraphrasing}\\RADAR~\cite{hu2023radar}\\Bias study~\cite{liang2023biased}};
  \node[blk,font=\scriptsize,text width=2.3cm,below] at (3.6,-0.75)
     {Binoculars~\cite{hans2024binoculars}\\Fast-DetectGPT~\cite{bao2024fast}\\Ghostbuster~\cite{verma2024ghostbuster}\\M4~\cite{wang2024m4}, MAGE~\cite{li2024mage}\\RAID~\cite{dugan2024raid}};
  \node[blk,font=\scriptsize,text width=1.8cm,above] at (4.9,0.45)
     {Survey~\cite{wu2025survey}};
\end{tikzpicture}
\caption{Timeline of representative work on AI-generated text detection}
\label{fig:timeline}
\end{figure}

% ---------------------------------------------------------------------
\section{Statistical and Heuristic Detection Methods}
% ---------------------------------------------------------------------

Statistical methods exploit the fact that a language model assigns
systematically different probabilities to its own output than to human
writing. Most of them need only a reference model to score tokens and either
no training at all (zero-shot) or a very small classifier on top of a few
features.

\subsection{GLTR: Token Rank and Probability}

GLTR (Giant Language model Test Room)~\cite{gehrmann2019gltr} was one of the
first practical tools. For each token of a passage it computes, under a
reference model such as GPT-2, the probability of the observed token, its
rank in the predicted distribution and the entropy of that distribution. It
then colours each token by its rank bucket (top-10, top-100, top-1000, other).
Generated text is dominated by top-10 tokens, while human text contains a
noticeable share of low-rank choices. In a user study, GLTR raised the rate at
which untrained people correctly identified generated text from roughly chance
level to over 70\%. GLTR established that simple distributional statistics
separate human and machine text, and its features form the basis of almost
every heuristic that followed, including the perplexity features used in this
project.

\subsection{DetectGPT and Fast-DetectGPT: Probability Curvature}

DetectGPT~\cite{mitchell2023detectgpt} proposed a zero-shot criterion based on
the local shape of the log-probability function. Text sampled from a model
tends to lie near a local maximum of that model's log-probability; small
rewrites of such text therefore lower its log-probability on average, while
rewrites of human text change it in either direction. DetectGPT generates
many perturbed versions of a passage with a mask-filling model (T5), scores
the original and the perturbations, and uses the normalised difference as its
statistic. The method needs no classifier training and performed strongly on
several generators, but it requires on the order of a hundred model calls per
document, which makes it expensive for bulk screening.

Fast-DetectGPT~\cite{bao2024fast} replaces the perturbation step with a
sampling step based on \emph{conditional} probability curvature, which can be
computed from a single forward pass. The authors report that it accelerates
detection by a factor of about 340 relative to DetectGPT while improving
detection accuracy by a relative margin of roughly 75\% in their experiments.
This shows that curvature-type statistics can be made cheap enough for the
first stage of a pipeline.

\subsection{Binoculars: Cross-Perplexity}

Binoculars~\cite{hans2024binoculars} addresses a known weakness of raw
perplexity: a human text on an unusual topic may have high perplexity simply
because its content is surprising, while a model asked to write about that
topic would also produce surprising content. Binoculars therefore divides the
perplexity of a passage under an ``observer'' model by a cross-perplexity
computed from how surprising one model's predictions are to a closely related
``performer'' model. The ratio normalises away the effect of the prompt and
topic. Without any training on ChatGPT data, Binoculars detected over 90\% of
generated samples from ChatGPT and other LLMs at a false positive rate of
0.01\% in the authors' evaluation. Its main practical cost is that it must
run two billion-parameter-scale language models for every document.

\subsection{Stylometric and Burstiness Signals}

Long before neural text generation, authorship attribution relied on
stylometric features: sentence-length distribution, vocabulary richness,
function-word frequencies and part-of-speech patterns. These features carry
over naturally to machine-text detection. The analysis accompanying the HC3
corpus~\cite{guo2023hc3} found that ChatGPT answers differ from human expert
answers in measurable ways such as vocabulary size, the distribution of
part-of-speech tags and the use of connective words, and that a simple
classifier on such statistics already achieves useful accuracy. Burstiness,
the variation of sentence length and per-sentence perplexity across a
document, is the signal popularised by the original release of GPTZero.
Stylometric features are inexpensive and interpretable, but they are also
the features most affected by genre and by the writer's proficiency, which
is a concern discussed in Section~\ref{sec:lit-fairness}.

\subsection{Ghostbuster: Learned Combinations of Probability Features}

Ghostbuster~\cite{verma2024ghostbuster} sits between zero-shot and supervised
methods. It passes a document through a series of weaker language models
(a unigram model, a trigram model and two GPT-3 variants), computes token
probabilities under each, and runs a structured search over arithmetic
combinations of these probability vectors to produce a compact set of
features. A linear classifier is then trained on the selected features. The
method does not need the generating model's own probabilities and the authors
report better generalisation across domains, prompts and generators than
single-model classifiers. Ghostbuster is the closest published precedent for
the heuristic stage of \veriscript{}: a small, interpretable classifier on top
of language-model-derived features.

Table~\ref{tab:lit-stat} summarises the statistical methods reviewed.

\begin{table}[H]
\caption{Summary of statistical and heuristic detection methods}
\label{tab:lit-stat}
\centering
\small
\begin{tabularx}{\textwidth}{|p{2.5cm}|L|L|L|}
\hline
\textbf{Method} & \textbf{Core signal} & \textbf{Strength} & \textbf{Limitation} \\ \hline
GLTR~\cite{gehrmann2019gltr} & Token rank, probability and entropy under a reference LM & Very fast, visual, interpretable & Shallow; weak against modern sampling and paraphrasing \\ \hline
DetectGPT~\cite{mitchell2023detectgpt} & Log-probability curvature under perturbations & Zero-shot, strong on many generators & About 100 model calls per document \\ \hline
Fast-DetectGPT~\cite{bao2024fast} & Conditional probability curvature & Single pass; much faster than DetectGPT & Still needs a sizeable scoring LM \\ \hline
Binoculars~\cite{hans2024binoculars} & Ratio of perplexity to cross-perplexity & High accuracy at very low FPR without training & Two large LMs per document \\ \hline
Stylometry / burstiness~\cite{guo2023hc3} & Sentence-length variance, TTR, POS patterns & Cheap, fully explainable & Sensitive to genre and writer proficiency \\ \hline
Ghostbuster~\cite{verma2024ghostbuster} & Searched combinations of weak-LM probabilities + linear model & Good cross-domain generalisation & Multiple LM passes; feature search is costly offline \\ \hline
\end{tabularx}
\end{table}

% ---------------------------------------------------------------------
\section{Transformer-Based Classifiers}
% ---------------------------------------------------------------------

The second major family treats detection as ordinary supervised text
classification and fine-tunes a pretrained encoder on labelled examples.

\subsection{Early Neural Detectors}

When OpenAI staged the release of GPT-2, it published a detector obtained by
fine-tuning RoBERTa on GPT-2 outputs~\cite{solaiman2019release,liu2019roberta}.
The fine-tuned classifier detected GPT-2 text far more accurately than
zero-shot baselines, and for some time it was the reference detector in the
field. At the same time, Zellers \textit{et al.}~\cite{zellers2019grover}
introduced Grover, a model for generating fake news articles, and showed that
the best defence against Grover was Grover itself used as a discriminator: a
detector benefits greatly from having been trained on the same kind of model
that produced the text. This observation foreshadows the generalisation
problem discussed later.

Ippolito \textit{et al.}~\cite{ippolito2020automatic} compared automatic
detectors with human judges and found an instructive asymmetry. Decoding
strategies that remove the low-probability tail (such as top-$k$) make text
\emph{more} convincing to humans but \emph{easier} for automatic detectors,
because the missing tail is a statistical fingerprint. Conversely, sampling
from the full distribution fools detectors more often but produces text that
humans find less fluent. Automatic and human detection therefore exploit
different evidence.

\subsection{Stronger Encoders: RoBERTa and DeBERTa}

Later work replaced the original classifier with stronger encoders.
RoBERTa~\cite{liu2019roberta} improves on BERT~\cite{devlin2019bert} by
training longer on more data with dynamic masking and without the
next-sentence objective. DeBERTa~\cite{he2021deberta} introduces disentangled
attention, in which each token is represented by separate content and
relative-position vectors, together with an enhanced mask decoder;
DeBERTa-v3~\cite{he2023debertav3} adopts an ELECTRA-style replaced-token
detection objective and is among the strongest base-sized encoders available.
The HC3 authors fine-tuned RoBERTa detectors on their corpus and obtained
very high in-domain scores~\cite{guo2023hc3}, and fine-tuned RoBERTa and
DeBERTa models are standard strong baselines in recent benchmarks~\cite{wang2024m4,
li2024mage,dugan2024raid}.

The weakness of these classifiers is not in-domain accuracy but cost and
transfer. A base-sized encoder has on the order of 125--185 million parameters
and processes up to 512 tokens per forward pass; long documents need several
passes. On a CPU-only server, inference for a large batch of essays is slow.
The benchmarks reviewed in Section~\ref{sec:lit-bench} also show that accuracy
falls when the test text comes from a generator, domain or decoding strategy
not seen in training.

\subsection{Adversarially Trained Detectors}

RADAR~\cite{hu2023radar} targets robustness to paraphrasing by training a
detector and a paraphraser together in an adversarial loop. The paraphraser
is rewarded for rewriting AI text so that the detector labels it human; the
detector is trained on human text, original AI text and the paraphraser's
output. Evaluated on eight LLMs across four datasets, RADAR outperformed
existing detectors, especially when paraphrasing was applied, and its
detector transferred well across instruction-tuned models. RADAR's detector is
itself a fine-tuned RoBERTa-large model, so it shares the inference cost of the
transformer family. Adversarial training is complementary to \veriscript{}:
a RADAR-style model could serve as the second stage of the cascade.

% ---------------------------------------------------------------------
\section{Commercial Market Solutions}
% ---------------------------------------------------------------------

Commercial detectors are what institutions actually deploy, so their
characteristics define the practical baseline for this project. Their models
are proprietary and their public documentation gives limited technical
detail; the descriptions below are therefore based on vendor descriptions and
on published independent evaluations, and are kept to properties that matter
for deployment.

\begin{itemize}
  \item \textbf{Turnitin (AI writing detection)} is integrated into the
        similarity-checking workflow that universities already use through
        their learning management systems (LMS). It reports the proportion
        of a submission it considers AI-generated. Its main advantages are
        seamless integration and a large academic customer base; its main
        disadvantages for this project's purposes are that it is a closed
        model sold through institutional licences and that the basis for an
        individual score cannot be inspected or audited by the institution.
  \item \textbf{GPTZero} was first released as a lightweight detector built on
        perplexity and burstiness and has since grown into a multi-stage
        product that combines statistical signals with learned models. It is
        fast and accessible, with free and paid tiers. As with other
        perplexity-based tools, formal and highly structured human prose is a
        known source of false alarms.
  \item \textbf{Copyleaks} offers AI-content detection through a web interface
        and an API, with support for several languages. Processing takes place
        on the vendor's servers and is priced per page or per credit, which
        adds network latency and recurring cost for bulk use.
  \item \textbf{Originality.ai} is marketed mainly to web publishers and
        search-engine-optimisation (SEO) professionals. It emphasises
        sensitivity to text from recent models; high sensitivity settings
        raise the risk of flagging polished human writing.
  \item \textbf{Scribbr and QuillBot} are writing-assistance platforms that
        include AI detectors alongside paraphrasing and grammar tools. Free
        tiers impose character limits, and the tools are designed for
        individual checks rather than bulk institutional screening.
\end{itemize}

Table~\ref{tab:lit-commercial} compares these products on the deployment
dimensions identified in Chapter~\ref{ch:intro}. The common pattern is that
every commercial tool is closed, none publishes the features behind its
decision, and cost scales with volume.

\begin{table}[H]
\caption{Deployment characteristics of commercial AI-text detectors}
\label{tab:lit-commercial}
\centering
\small
\begin{tabularx}{\textwidth}{|p{2.6cm}|L|L|L|L|}
\hline
\textbf{Tool} & \textbf{Primary audience} & \textbf{Access model} & \textbf{Explains decisions?} & \textbf{Main concern for institutions} \\ \hline
Turnitin & Universities (LMS) & Institutional licence & Score only & Opaque, cannot be audited \\ \hline
GPTZero & Educators, individuals & Free and paid tiers, API & Sentence highlighting & False alarms on formal prose \\ \hline
Copyleaks & Enterprises, education & Per-page / credit API & Score with highlights & Recurring cost, server-side latency \\ \hline
Originality.ai & Publishers, SEO & Pay-per-scan credits & Score with highlights & Aggressive thresholds, false positives \\ \hline
Scribbr / QuillBot & Students, writers & Freemium, character limits & Score only & Not suited to bulk screening \\ \hline
\end{tabularx}
\end{table}

% ---------------------------------------------------------------------
\section{Watermarking and Generation-Time Signals}
% ---------------------------------------------------------------------

Watermarking takes a structurally different approach: instead of detecting
text after the fact, the model provider embeds a detectable signal at
generation time. Kirchenbauer \textit{et al.}~\cite{kirchenbauer2023watermark}
proposed seeding a pseudo-random function with the previous token to split
the vocabulary into a ``green'' list and a ``red'' list at every step, and
adding a small bias to the logits of green tokens. Human text contains green
tokens about as often as chance predicts, while watermarked text contains
significantly more; a one-sided $z$-test on the green-token count detects the
watermark from as few as a couple of hundred tokens, with negligible effect on
text quality.

Watermarking is extremely cheap to detect and comes with a statistical
guarantee, but it has two limitations that rule it out as the primary method
for this project. First, it requires the cooperation of the model provider;
it cannot be applied retroactively to text produced by an unwatermarked model,
an open-weight model run locally, or a provider that chooses not to
participate. Second, paraphrasing substantially weakens the
watermark~\cite{krishna2023paraphrasing,sadasivan2023reliably}. Watermark
checking is therefore treated in this report as a possible future third tier
of the cascade rather than as an alternative to it.

% ---------------------------------------------------------------------
\section{Benchmark Datasets and Generalisation}
\label{sec:lit-bench}
% ---------------------------------------------------------------------

Detector accuracy reported on a single dataset says little about behaviour in
deployment. Several benchmarks have been built specifically to measure
generalisation.

\textbf{HC3}~\cite{guo2023hc3} (Human ChatGPT Comparison Corpus) pairs human
and ChatGPT answers to the same questions. Its English portion contains 24,322
questions with 58,546 human answers and 26,903 ChatGPT answers drawn from
Reddit's ELI5, open-domain question answering, Wikipedia computer-science
topics, finance and medicine. Because both answers respond to the same prompt,
HC3 allows a controlled comparison in which topic is held fixed.

\textbf{M4}~\cite{wang2024m4} is a multi-generator, multi-domain and
multilingual benchmark covering several generators (including ChatGPT,
GPT-3.5 davinci, Cohere, Dolly-v2 and BLOOMz), domains such as Wikipedia,
WikiHow, Reddit, arXiv abstracts and peer reviews, and several languages. Its
main finding is that detectors struggle to generalise to unseen domains or
unseen generators, and that in such cases they tend to misclassify
machine-generated text as human-written.

\textbf{MAGE}~\cite{li2024mage} collects 447,674 human-written and
machine-generated texts from ten writing tasks and 27 LLMs, organised into
testbeds of increasing difficulty. Detection becomes markedly harder
out of distribution; the best detector identified 84.12\% of out-of-domain
texts generated by a new LLM, which the authors interpret as showing that the
problem is difficult but feasible.

\textbf{RAID}~\cite{dugan2024raid} is the largest of these benchmarks, with
over six million generations from 11 models across 8 domains, several decoding
strategies and 11 adversarial attacks such as paraphrasing, synonym swaps and
homoglyph substitution. RAID showed that detectors that claim very high
accuracy on their own test sets lose much of that accuracy when evaluated on
unseen models, on different sampling settings, or under simple attacks.

Table~\ref{tab:lit-bench} summarises the benchmarks. Together they motivate two
design decisions in \veriscript{}: the training corpus must mix generators and
domains, and the evaluation must include a held-out generator.

\begin{table}[H]
\caption{Benchmark datasets for machine-generated text detection}
\label{tab:lit-bench}
\centering
\small
\begin{tabularx}{\textwidth}{|p{2.0cm}|L|L|L|}
\hline
\textbf{Benchmark} & \textbf{Composition} & \textbf{Key finding} & \textbf{Use in this project} \\ \hline
HC3~\cite{guo2023hc3} & 24,322 English questions; 58,546 human and 26,903 ChatGPT answers; 5 domains & Human and ChatGPT answers differ in vocabulary and style & Topic-controlled training and validation pairs \\ \hline
M4~\cite{wang2024m4} & Multiple generators, domains and languages & Poor transfer to unseen domains and generators & Domain diversity; English subset \\ \hline
MAGE~\cite{li2024mage} & 447,674 texts; 27 LLMs; 10 tasks & Out-of-distribution detection is much harder & Additional out-of-distribution test material \\ \hline
RAID~\cite{dugan2024raid} & 6M+ generations; 11 models; 8 domains; 11 attacks & High reported accuracy collapses on unseen models and attacks & Robustness and attack evaluation \\ \hline
\end{tabularx}
\end{table}

% ---------------------------------------------------------------------
\section{Robustness, Reliability and Fairness}
\label{sec:lit-fairness}
% ---------------------------------------------------------------------

\subsection{Paraphrasing Attacks}

Krishna \textit{et al.}~\cite{krishna2023paraphrasing} trained DIPPER, an
11-billion-parameter paraphrase model that rewrites whole paragraphs with
controllable lexical diversity and reordering. Paraphrasing text generated by
three LLMs with DIPPER evaded several detectors, including watermarking,
GPTZero, DetectGPT and OpenAI's text classifier; for example, it reduced
DetectGPT's detection accuracy from 70.3\% to 4.6\% at a fixed 1\% false
positive rate without appreciably changing the meaning. The authors proposed a
retrieval-based defence in which the model provider stores its generations and
checks candidate texts for semantic similarity, which again requires provider
cooperation.

\subsection{Theoretical Limits of Detection}

Sadasivan \textit{et al.}~\cite{sadasivan2023reliably} argued that as language
models approximate the distribution of human text more closely, the best
achievable detector performance approaches that of a random classifier, and
demonstrated recursive paraphrasing attacks against several detector types.
Their analysis does not make detection useless in practice, because current
models still differ measurably from human writing, but it is an important
caution: any detector's output is evidence, not proof, and should be presented
with a calibrated confidence rather than as a verdict.

\subsection{Bias Against Non-Native English Writers}

Liang \textit{et al.}~\cite{liang2023biased} evaluated seven widely used GPT
detectors on essays written by native and non-native English speakers. The
detectors classified native-written essays accurately, but more than half of
the TOEFL essays written by non-native speakers were incorrectly flagged as
AI-generated. The authors traced the effect to perplexity: writing with a more
limited vocabulary and simpler constructions is more predictable to a language
model, which is exactly the property perplexity-based detectors rely on. They
also showed that prompting an LLM to enrich the vocabulary of the non-native
essays reduced the false-positive rate, which demonstrates that the detectors
were responding to linguistic sophistication rather than authorship.

This finding is directly relevant to \veriscript{}, because its first stage
uses perplexity and lexical diversity. It motivates two design choices
explained in Chapter~\ref{ch:proposed}: the heuristic stage is allowed to
decide only when it is highly confident, so that borderline texts, which are
where such bias is most likely to cause errors, are passed to the
context-aware transformer; and the false-positive rate on non-native writing
is measured as a primary evaluation criterion.

% ---------------------------------------------------------------------
\section{Efficient Inference: Cascades and Early Exit}
% ---------------------------------------------------------------------

The cost of transformer inference is a general problem, and several lines of
work outside the detection literature address it. They provide the
methodological basis for the cascade in \veriscript{}.

\textbf{Boosted cascades.} The Viola--Jones face detector~\cite{viola2001rapid}
arranged increasingly complex classifiers in a cascade so that the great
majority of image windows were rejected by the first, very cheap stages.
This made real-time detection possible on 2001 hardware and is the classic
demonstration that most inputs do not need the most expensive model.

\textbf{Early exit in deep networks.} BranchyNet~\cite{teerapittayanon2016branchynet}
attached side classifiers to intermediate layers of a convolutional network
and let a sample exit at the first branch whose prediction entropy fell below
a threshold. DeeBERT~\cite{xin2020deebert} applied the same idea to BERT,
inserting an exit after each transformer layer and reporting substantial
inference savings with little loss in accuracy on standard classification
tasks.

\textbf{Model compression.} DistilBERT~\cite{sanh2019distilbert} takes the
complementary approach of training a smaller student network to imitate a
larger teacher, retaining most of BERT's language-understanding performance
with 40\% fewer parameters and 60\% faster inference.

\textbf{Cascades of language-model services.} FrugalGPT~\cite{chen2023frugalgpt}
cascades LLM APIs of increasing price, accepting a cheaper model's answer when
a learned scoring function judges it reliable. The authors report matching the
strongest model's quality at a large reduction in cost.

These works establish that confidence-based routing can preserve most of an
expensive model's accuracy while avoiding most of its cost, provided the cheap
stage's confidence is well calibrated. None of them, however, addresses
AI-text detection, and none combines an \emph{interpretable feature-based}
first stage with a transformer second stage.

% ---------------------------------------------------------------------
\section{Synthesis and Research Gap}
% ---------------------------------------------------------------------

Table~\ref{tab:lit-summary} brings together the representative approaches
reviewed in this chapter. The accuracy and latency ratings are qualitative and
relative to one another, reflecting the evidence reported in the cited works.

\begin{table}[H]
\caption{Summary of representative detection approaches and market tools}
\label{tab:lit-summary}
\centering
\small
\begin{tabularx}{\textwidth}{|p{2.9cm}|p{2.3cm}|p{1.65cm}|p{1.65cm}|L|}
\hline
\textbf{Approach / Tool} & \textbf{Category} & \textbf{Accuracy} & \textbf{Latency} & \textbf{Key limitation} \\ \hline
GLTR~\cite{gehrmann2019gltr} & Heuristic & Low--Mod. & Very low & Shallow feature depth \\ \hline
DetectGPT~\cite{mitchell2023detectgpt} & Zero-shot curvature & Moderate & High & Many model calls per sample \\ \hline
Fast-DetectGPT~\cite{bao2024fast} & Zero-shot curvature & High & Low--Mod. & Needs a large scoring LM \\ \hline
Binoculars~\cite{hans2024binoculars} & Cross-perplexity & High & Moderate & Requires two LLMs \\ \hline
Ghostbuster~\cite{verma2024ghostbuster} & Feature-based + linear & High & Moderate & Several LM passes \\ \hline
RoBERTa / DeBERTa~\cite{liu2019roberta,he2023debertav3} & Fine-tuned transformer & Very high (in-domain) & High & GPU-dependent; weak transfer \\ \hline
RADAR~\cite{hu2023radar} & Adversarial transformer & High & High & Large encoder; costly training \\ \hline
Watermark~\cite{kirchenbauer2023watermark} & Generation-time & Very high & Very low & Needs provider cooperation \\ \hline
Turnitin & Commercial LMS & High & Moderate & Black box, licence cost \\ \hline
GPTZero & Commercial web & Mod.--High & Low--Mod. & False positives on formal text \\ \hline
Copyleaks & Commercial API & High & Moderate & Cost per page \\ \hline
Originality.ai & Commercial web & High (sensitive) & Moderate & High false-positive risk \\ \hline
\textbf{VeriScript (this work)} & \textbf{Hybrid cascade} & \textbf{High (F1 0.971 on test)} & \textbf{Moderate (GPT-2 first stage)} & \textbf{First stage as costly as the second (Ch.~7)} \\ \hline
\end{tabularx}
\end{table}

Five observations emerge from the review.
\begin{enumerate}
  \item \textbf{Methods are evaluated in isolation.} Academic work reports the
        accuracy of individual detectors but rarely the accuracy--latency
        trade-off that determines whether a detector can be deployed at
        institutional scale.
  \item \textbf{Cheap and accurate detectors are not combined.} Feature-based
        detectors and fine-tuned transformers are presented as alternatives.
        Although cascades and early exit are well established for reducing
        inference cost in other fields~\cite{viola2001rapid,xin2020deebert,
        chen2023frugalgpt}, no surveyed detection system routes documents
        between an interpretable heuristic stage and a transformer stage on the
        basis of calibrated confidence.
  \item \textbf{Commercial tools are closed and costly.} The tools actually
        used by institutions do not explain their decisions and charge per
        licence or per page.
  \item \textbf{Generalisation is weak.} Benchmarks consistently show large
        accuracy drops on unseen generators, domains and attacks~\cite{wang2024m4,
        li2024mage,dugan2024raid,krishna2023paraphrasing}.
  \item \textbf{Fairness is under-evaluated.} Perplexity-based detection is
        biased against non-native writers~\cite{liang2023biased}, yet few
        systems measure false-positive disparity as part of their evaluation.
\end{enumerate}

\veriscript{} is designed to address these gaps together: an open,
cost-aware, interpretable cascade that is evaluated jointly on accuracy,
efficiency, cross-generator robustness and fairness. Chapter~\ref{ch:problem}
turns these gaps into a formal problem statement and a set of objectives.
